\section{Sim2Real 与策略蒸馏}

\subsection{Sim2Real 的分布转移策略}

将实验室训练迁移到物理世界时，Dynamics gap 会显著降低性能。常见对策包括：domain randomization（随机化环境参数）、system identification（模型化真实动力学）、解析 backup policy 作为 fallback。lecture 中强调的 `Sim-to-Real index` 通过对比 sim rollout 和 real rollout 的 total variation distance 得出 policy drift 指标。

\subsection{策略蒸馏与小模型部署}

Distillation pipeline 一般分为：teacher run -> dataset harvest -> student training -> student evaluation。优质的 dataset 需要覆盖 teacher 出错的 edge cases。student training 中通常使用 `mse` + `kl` loss 复合 objective 维持行为一致性。

\begin{knowledgebox}{为何 distillation 还是必要}
即使 teacher 能取得最尖端表现，因其体量大、推理成本高，仍需 distill 到 latency-friendly policy。蒸馏过程将 teacher 的 knowledge 以 soft target 形式写入 student，使 student 可以在 resource-constrained device 上运行。
\end{knowledgebox}

\subsection{本章小结}

Sim2Real 需要 domain randomization、system ID 与 fallback policies。策略蒸馏将 teacher 的能力转移到更轻量 student，是部署到移动/嵌入式平台的关键途径。


